\chapter{Metodologia y herramientas de seguimiento}

\section{Metodologia}
El desarrollo se aborda mediante una metodología iterativa e incremental. Esta elección responde a la naturaleza experimental del proyecto: integrar IA local, contenedores, domótica, RAG y acceso remoto seguro implica incertidumbres técnicas que no pueden resolverse solo con diseño previo. Cada iteración debe producir un bloque funcional, medible y conectable al resto del ecosistema.

La construcción se organiza en sprints técnicos: infraestructura base, inferencia local, orquestación, integración sensorial, memoria documental, seguridad, monitorización y expansión de servicios. Al final de cada ciclo se revisan resultados, problemas, consumo de recursos, estabilidad y utilidad real.

\section{Herramientas de seguimiento}
El seguimiento del trabajo se apoya en control de versiones, documentación técnica y registro de decisiones. Las configuraciones de contenedores, scripts, flujos de automatización y notas de despliegue deben mantenerse versionadas para poder reproducir cambios y revertir errores. La memoria funcionará además como bitácora técnica: cada herramienta relevante tendrá una explicación de uso, motivación, configuración, problemas y papel dentro de Odín.

\section{Mecanismos de rigor y control}
El rigor se garantiza mediante validaciones por capas. Primero se comprueba la estabilidad de la infraestructura. Después se validan servicios individuales. Finalmente se prueban flujos completos: entrada, procesamiento, decisión, acción y persistencia. Para cada subsistema se documentarán criterios de aceptación, errores encontrados y medidas de mitigación.

\section{Evolución y endurecimiento de la metodología}
La metodología inicial se volvió más rigurosa a medida que aumentó el número de servicios y dependencias. Cada iteración terminó siguiendo un ciclo estable:

\begin{enumerate}
  \item formular una hipótesis o necesidad concreta;
  \item aplicar un cambio pequeño y reversible;
  \item comprobar el servicio de forma aislada mediante API, registro o prueba directa;
  \item validar después el flujo completo desde la interfaz utilizada por el usuario;
  \item documentar el resultado, las credenciales implicadas y el procedimiento de recuperación.
\end{enumerate}

Este cambio introdujo prácticas que no estaban suficientemente desarrolladas en el planteamiento inicial: configuraciones sanitizadas, comprobaciones de salud de contenedores, separación entre prueba técnica y aceptación funcional, inventario explícito de puertos y nombres locales y preferencia por modificaciones reversibles. También se adoptó como criterio que una respuesta del modelo no demuestra por sí sola que una acción se haya ejecutado; el resultado debe verificarse en el sistema de destino.

La consecuencia principal fue un desplazamiento de esfuerzo desde la instalación hacia la integración, el diagnóstico y la documentación reproducible. Este ajuste explica parte de la desviación respecto a la planificación original y constituye uno de los aprendizajes centrales del proyecto.

\section{Profundización en herramientas}
Una prioridad de la memoria final será explicar con detalle las herramientas utilizadas, no como una lista decorativa sino como decisiones de ingeniería. Para cada herramienta principal se incluirá:
\begin{itemize}
  \item problema concreto que resuelve dentro de Odín;
  \item motivo de elección frente a alternativas;
  \item arquitectura de despliegue y configuración relevante;
  \item integración con el resto de servicios;
  \item limitaciones encontradas durante el desarrollo;
  \item evidencias de funcionamiento mediante capturas, logs, tablas o pruebas.
\end{itemize}
